\scanpage{81}
\mspage{25}

Cela prouve déjà que les foncteurs $r$ et $C$ induisent des équivalences \unsure{quasi-inverses} l'une de l'autre sur les catégories localisées. Par le lemme p.~11, cela montre que pour tte sous-catégorie pleine \inserted{$\unsure{\mathcal F|\mathscr F|F}$} de $\mathrm{Cat}_{/S}$ contenant $\underline{\mathrm{Parf}}\,S$, l'inclusion de $\mathcal F$ dans $\mathrm{Cat}_{/S}$ induit une équivalence sur les localisés, en prenant comme localiseur dans $\mathcal F$ le localiseur ``grossier'' induit par $\mathrm{Cat}_{/S}$. Ceci s'applique notamment aux sous-catégories pleines
\[
\underline{\mathrm{Liss}}\,S,\quad\underline{\mathrm{Fib\,hot}}\,S,\quad\underline{\mathrm{Prop}}\,S,
\]
% vertical stroke beside the following four manuscript lines
\leavevmode\llap{\smash{\rule[-2.3\baselineskip]{.4pt}{3\baselineskip}}\quad}mais en prenant dans $\underline{\mathrm{Liss}}\,S$ et $\underline{\mathrm{Prop}}\,S$ \unsure{n'ont|non} pas le localiseur fin ``fibre par fibre'', comme je l'avais annoncé à tort, mais le localiseur grossier induit par $\mathrm{Cat}_{/S}$. (Pour $\underline{\mathrm{Fib\,hot}}\,S$, il n'y a pas de risque d'ambiguïté). Il faudrait que je me \unsure{convainque} que si je prends là les $W$-équivalences fines (dans le cas $W=W_{\infty}$, disons), je ne trouve \uline{pas} les ``bons'' localisés -- i.e. le foncteur\unsure{,|.} des localisés, vers le localisé de $\mathrm{Cat}_{/S}$, n'est pas une équivalence.

Ceci me fait \unsure{supposer} que dans $\underline{\mathrm{Fib}}\,S$ et $\underline{\mathrm{Cofib}}\,S$ également, c'est \deleted{\unsure{le|la}}l'\deleted{\unsure{notion|relation}}\inserted{équivalence} \uline{grossière}

% archive sheet number: 80
